\section{Physical intrinsic resource certification}
\label{sec:v15-certification}
We can now state a single conclusion joining the analytic and finite
parts of the theory. The original experiments, not their approximations,
are the objects whose intrinsic deficiency is bounded. Fix one of the
private, hidden-selector, or visible-selector signatures, denoted $b$;
its state widths and, where relevant, selector cardinality are fixed.
Write $\delta_b(E,F)$ for the corresponding intrinsic deficiency, with
source $E$ simulating target $F$. The marked feedback comparison always
uses the actual initial mark. For a visible selector its reference law
is the joint law specified in Section~\ref{sec:v13-deficiency}.

\begin{theorem}[Same-budget physical certification]
\label{thm:v15-main}
Let $E,F$ be finite-horizon active Gaussian acquisitions on Hilbert
probability models of the form considered above, with finite parameter,
action and actual-mark alphabets, positive noise variances, and certified
$L^2$ preparation bounds. For each model choose a finite trial
presentation with rigorous marked error $\alpha_E,\alpha_F$, supplied
either by Theorem~\ref{thm:v14-residual} or by
Theorem~\ref{thm:v15-enriched}. Choose finite report channels as in
Lemma~\ref{lem:v15-quantization} and rational marked trees as in
Theorem~\ref{thm:v15-rational}, with certified errors $\beta_E,\beta_F$
and $\rho_E,\rho_F$. Put
\[
 a=1\wedge(\alpha_E+\beta_E+\rho_E),\qquad
 c=1\wedge(\alpha_F+\beta_F+\rho_F).
\]
Let $[L,U]$ be any verified local-certificate interval for
$\delta_b(E_{\mathbb Q},F_{\mathbb Q})$ at the original signature $b$.
Then the original physical deficiency belongs to
\begin{equation}\label{eq:v15-main-interval}
 \left[\max\{0,L-a-c\},\ \min\{1,U+a+c\}\right].
\end{equation}
An upper witness is an executable simulator of the original physical
experiment with exactly the declared persistent-state and selector
budget. A strict lower endpoint is an impossibility certificate at that
same budget; it does not refer to a convexified private feasible set.

For the effective assertion, the noise scales are computable reals
with supplied rational positive lower bounds. Under effective graph
data, including the marked cell-integral
operation in Definition~\ref{def:v15-effective}, there is a terminating
procedure producing intervals of the form \eqref{eq:v15-main-interval}
of arbitrarily small prescribed width. Their nested intersections
converge to $\delta_b(E,F)$. The procedure uses only finite rational
certificates and certified analytic input operations. It does not assume
attainment of the original nonfinite optimization.
\end{theorem}
\begin{proof}
Each of the three presentation changes is implemented by a stateless,
seed-free, parameter-blind and task-blind report channel. The physical
trial change uses the identity report channel; the quantization change
uses $Q$ in one direction and transiently randomized $H$ in the other;
the rationalization change uses the identity finite-report channel.
Their compositions give two-sided marked feedback comparisons
$E\leftrightarrow E_{\mathbb Q}$ with error $a$ and
$F\leftrightarrow F_{\mathbb Q}$ with error $c$. Fresh per-report
randomness in a presentation channel is integrated into a simulator row;
it is not a persistent public or hidden mode. All comparisons include
the joint reference of a visible selector.

For an executable finite upper witness, first pass a report of $E$
through the forward presentation, apply the finite simulator, and pass
its output through the reverse presentation for $F$. The stored state
and persistent selector are exactly those of the finite witness. The
marked feedback triangle inequality bounds its error by $U+a+c$.
Conversely, any physical simulator with error smaller than $L-a-c$
could be composed with the opposite stateless presentations to give a
finite simulator of the same signature with error smaller than $L$.
This contradicts the verified finite lower certificate. Equivalently,
Theorem~\ref{thm:v14-stability} gives
\[
 |\delta_b(E,F)-\delta_b(E_{\mathbb Q},F_{\mathbb Q})|\le a+c.
\]
Taking infima in this argument requires no nonfinite minimizer.

For effectiveness, first choose enriched tolerances so that the two
rational upper bounds for $\alpha$ are as small as requested. The
terminating search of Theorem~\ref{thm:v15-enriched} supplies the trial
means and their norm bounds. Choose rational $R$ large and $h$ small
until the explicit rational upper bounds for $\beta$ are small. There
are then finitely many marked prefix masses. Request their enclosures
at a precision making the finite sum in \eqref{eq:v15-rational} small,
and normalize rational lower endpoints. This gives actual finite
rational instruments, not floating-point transition tables.

For these fixed finite instruments, Theorem~\ref{thm:v14-local} supplies
convergent dyadic local lower certificates and executable rational
upper witnesses. Refine until $U-L$ is below the requested remaining
budget; convergence, or the explicit causal oscillation modulus, ensures
termination. Although the finite alphabet and row dimension may grow
as the physical tolerance decreases, at each stage this is a finite
problem and the mesh is selected for that problem. Thus the total
interval width is at most $U-L+2(a+c)$ and can be made arbitrarily
small. Intersect each interval with its predecessors. They all contain
the same true value, so the intersections are nonempty and nested,
and their lengths tend to zero.
\end{proof}

\begin{remark}[What is constructive in the theorem]
The quantitative implication \eqref{eq:v15-main-interval} holds for
any rigorously supplied analytic enclosures. Its effective-convergence
assertion has an explicit analytic data hypothesis. A graph core known
only to exist is not, by itself, a numerical integration algorithm.
Neither arbitrary real preparations nor arbitrary measurable marks
are silently declared computable. Section~\ref{sec:v15-physical-data}
gives genuine hard-sphere data for which the required calculations can
be certified, including an exact continuous-report deficiency.
The private certificate search remains nonconvex and may be expensive;
the witness-dimension theorem reduces the number of distinct tests,
not the number of cells or the cost of real quantifier elimination.
\end{remark}

\begin{corollary}[Regenerative risk certificates]
\label{cor:v15-regenerative}
Retain the actual synchronous reset and stationary-policy signature of
Section~\ref{sec:v13-regeneration}. For losses in $[0,1]$, certified
marked presentation errors $\epsilon_j$ for ages $0\le j\le N$ give
a uniform stationary-risk error at most
\[
 \eta\sum_{j=0}^N(1-\eta)^j\epsilon_j+(1-\eta)^{N+1}.
\]
The same bound holds for each lower support function of the stationary
risk body. Finite-age rational local certificates may be transferred
at their declared budgets by Theorem~\ref{thm:v15-main}. Under its
effective hypotheses the presentation part of the displayed risk
bound can be made arbitrarily small, and then $N$ can be increased.
\end{corollary}
\begin{proof}
Condition the stationary regenerative law on its age. The law of age
$j$ has weight $\eta(1-\eta)^j$, and the omitted ages have total weight
$(1-\eta)^{N+1}$. Apply the marked bounded-loss inequality to each
retained age and then sum. The bound precedes optimization, so it also
bounds lower support functions. Finite-age certificates concern those
finite-age experiments; this corollary does not identify a finite-age
optimized deficiency with an unrestricted infinite-history deficiency.
The register is actually cleared at reset and any persistent selector
remains charged, exactly as in the inherited theorem.
\end{proof}
